\documentclass[10pt,twoside]{article}
\usepackage[utf8]{inputenc}
\usepackage{amsmath}
\usepackage{amsfonts}
\usepackage{amssymb}
\usepackage[spanish,es-noshorthands]{babel}
\usepackage[T1]{fontenc}
\usepackage{lmodern}
\usepackage{graphicx,hyperref}
\usepackage{tikz,pgf}
\usepackage{multicol}
\usepackage{subfig}
\usepackage[papersize={6.5in,8.5in},width=5.5in,height=7in]{geometry}
\usepackage{fancyhdr}
\pagestyle{fancy}
\fancyhead[LE]{\includegraphics[height=12pt]{Images/logo-colegio.png} Aritmética $6^{\circ}$}
\fancyhead[RE]{}
\fancyhead[RO]{\textit{Germ\'an Avenda\~no Ram\'irez, Lic. U.D., M.Sc. U.N.}}
\fancyhead[LO]{}

\author{Germ\'an Avenda\~no Ram\'irez, Lic. U.D., M.Sc. U.N.}
\title{\begin{minipage}{.2\textwidth}
\includegraphics[height=1.75cm]{Images/logo-colegio.png}\end{minipage}
\begin{minipage}{.55\textwidth}
\begin{center}
Taller 13, Fracciones   \\
Aritmética $6^{\circ}$
\end{center}
\end{minipage}\hfill
\begin{minipage}{.2\textwidth}
\includegraphics[height=1.75cm]{Images/logo-sed.png} 
\end{minipage}}
\date{}
\begin{document}
\maketitle
Nombre: \hrulefill Curso: \underline{\hspace*{44pt}} Fecha: \underline{\hspace*{2.5cm}}
\section*{Nivel I}
\begin{enumerate}
 \item Enumera los términos de una fracción y di qué indica cada uno de ellos. Pon varios ejemplos.
\item ¿Qué fracción de hora son 20 minutos? Y ¿35 minutos? Y ¿55 minutos?
\item Para elaborar un tarro de frutas se han necesitado 400 gramos de plátanos, 350 gramos de fresas, 250 gramos de azúcar y 50 gramos de manzanas. ¿Qué fracción del total representa cada uno de estos productos?
\item Calcula
\begin{enumerate}\begin{multicols}{4}
 \item $\dfrac{5}{10}$ de 90
 \item $\dfrac{7}{9}$ de 72
 \item $\dfrac{4}{3}$ de 42
 \item $\dfrac{5}{9}$ de 540
\end{multicols}
 \end{enumerate}
\item  En una clase de 24 alumnos 5/8 son chicos. ¿Cuántos chicos y chicas hay en clase?
\item El depósito de un coche tiene una capacidad de 63 litros de gasolina, si gasta los 5/9 en una excursión, ¿cuántos litros le quedan al volver de viaje?
\item En la puerta de un cine hay 12 mujeres por cada 8 hombres y 16 niños. ¿Cuál es la relación entre hombres y mujeres? ¿Entre hombres y niños? Y ¿Entre mujeres y niños?
\item Indica de las siguientes fracciones cuáles dan como resultado un número natural y cuáles un número decimal:

3/2, 12/3, 17/4, 27/9, 14/6, 19/8, 21/7
\item Indica en las fracciones siguientes cuáles son mayores, iguales o menores que la unidad:

1/7, 3/5, 9/9, 5/3, 17/2, 2/9, 16/4, 18/17, 5/5, 6/9, 12/12.
\item Las fracciones siguientes son menores que la unidad. ¿Qué fracción falta en cada una de ellas para completar la unidad?

3/7, 3/8, 9/12, 11/16, 7/13, 5/9.
\item ¿Qué fracción sobra en cada una de las siguientes para obte-
ner la unidad?

6/5, 5/4, 16/9, 8/6, 17/13, 25/19.
\item ¿Entre que números naturales consecutivos están comprendidas las fracciones siguientes? 

7/5, 12/5, 12/3, 18/7, 21/5 .
\item Representa en la recta numérica las siguientes fracciones:

3/9, 4/5, 6/8, 4/4, 5/3, 6/4, 8/2.
\item Escribe dos fracciones equivalentes a cada una de estas:

13/5, 7/14, 5/2, 45/18, 3/21.

Explica como lo has hecho.
\item  Son equivalentes las parejas de fracciones siguientes:
\begin{enumerate}\begin{multicols}{3}
 \item $\dfrac{15}{4}$ y $\dfrac{75}{35}$
 \item $\dfrac{33}{42}$ y $\dfrac{132}{168}$
 \item $\dfrac{17}{62}$ y $\dfrac{51}{185}$\end{multicols}
\end{enumerate}
\item  En un campeonato de atletismo uno de los saltadores de altura consigue saltar más de dos metros 13 veces de 52 intentos, su contrincante salta más de 2 metros 11 veces de 44 intentos. ¿Cuál de los dos ha ganado?
\item  Ordena de mayor a menor los siguientes grupos de fracciones y explica como lo has hecho:
\begin{enumerate}\begin{multicols}{2}
 \item $\dfrac{9}{8}$, $\dfrac{7}{8}$, $\dfrac{3}{8}$, $\dfrac{17}{8}$, $\dfrac{1}{8}$
 \item $\dfrac{6}{9}$, $\dfrac{6}{14}$, $\dfrac{6}{7}$, $\dfrac{6}{11}$, $\dfrac{6}{8}$
 \item $\dfrac{3}{6}$, $\dfrac{3}{4}$, $\dfrac{5}{10}$
 \item $\dfrac{8}{4}$, $\dfrac{4}{1}$, $\dfrac{6}{3}$
 \end{multicols}
\end{enumerate}
\item Realiza las siguientes operaciones simplificando los resultados.
\begin{enumerate}\begin{multicols}{3}
 \item $\dfrac{3}{4}+\dfrac{5}{6}+\dfrac{6}{8}$
 \item $\dfrac{6}{10}-\dfrac{2}{8}=$
 \item $\dfrac{7}{8}+\dfrac{3}{6}-\dfrac{5}{12}=$
 \end{multicols}
\end{enumerate}
\item Realiza las siguientes operaciones simplificando los resultados.
\begin{enumerate}
\begin{multicols}{3}
 \item $\dfrac{5}{8}\cdot \dfrac{3}{9} \cdot \dfrac{4}{6}=$
 \item $$
 \end{multicols}
\end{enumerate}

\end{enumerate}

\end{document}
